% GENERATED FILE. Edit canonical lessons, then run npm run generate:books.
% canonical-source-hash: fnv1a64:a879f6f84ba13954
% canonical-lessons: TA-C275-iraku, TA-C275-val, TA-C275-kutu, TA-C275-ceti, TA-C275-koti2

\chapter{Feather, Tail, Nest, Plant, Creeper}
\label{ch:ta-feather-tail-nest-plant-creeper}
\hlchaptermodality{275}

\begin{chapteropening}
\textbf{By the end of this chapter:} \emph{I can say a feather, a tail, a nest, a plant and a creeper.}

Complete the last lesson of chapter 275: I can say a feather, a tail, a nest, a plant and a creeper.
\end{chapteropening}

\section[\texorpdfstring{\ta{இறகு} (iṟaku)}{இறகு (iṟaku)}]{\ta{இறகு} (iṟaku) — a feather}

\label{lesson:TA-C275-iraku}



\begin{quote}
Before the new one: say the Tamil for God, then the Tamil for an animal.
\end{quote}

\subsection*{You'll want to know: \ta{இறகு}}
\textbf{\ta{இறகு}} — \emph{iṟaku} — \textquotedblleft{}a feather\textquotedblright{}.

\begin{grammarlens}[title={What you've built}]
One more word for this chapter.
\end{grammarlens}

\subsection*{Your turn}
\begin{itemize}
\raggedright
  \item \emph{Say it:} \emph{iṟaku}
  \item \emph{Say it:} \emph{iṟaku}, once more
  \item \emph{Say it:} say \emph{iṟaku}
  \item \emph{Recall:} say the Tamil for God, then the Tamil for an animal, then say \emph{iṟaku} again
\end{itemize}

\begin{tcolorbox}[breakable,colback=teal!4,colframe=teal!35!black,title={Before you move on}]
What does \ta{இறகு} mean? (\textquotedblleft{}A feather\textquotedblright{}.) Say it once more.
\end{tcolorbox}
\section[\texorpdfstring{\ta{வால்} (vāl)}{வால் (vāl)}]{\ta{வால்} (vāl) — a tail}

\label{lesson:TA-C275-val}



\begin{quote}
Before the new one: say the Tamil for an animal, then the Tamil for a feather.
\end{quote}

\subsection*{You'll want to know: \ta{வால்}}
\textbf{\ta{வால்}} — \emph{vāl} — \textquotedblleft{}a tail\textquotedblright{}.

\begin{grammarlens}[title={What you've built}]
One more word for this chapter.
\end{grammarlens}

\subsection*{Your turn}
\begin{itemize}
\raggedright
  \item \emph{Say it:} \emph{vāl}
  \item \emph{Say it:} \emph{vāl}, once more
  \item \emph{Say it:} say \emph{vāl}
  \item \emph{Recall:} say the Tamil for an animal, then the Tamil for a feather, then say \emph{vāl} again
\end{itemize}

\begin{tcolorbox}[breakable,colback=teal!4,colframe=teal!35!black,title={Before you move on}]
What does \ta{வால்} mean? (\textquotedblleft{}A tail\textquotedblright{}.) Say it once more.
\end{tcolorbox}
\section[\texorpdfstring{\ta{கூடு} (kūṭu)}{கூடு (kūṭu)}]{\ta{கூடு} (kūṭu) — a nest}

\label{lesson:TA-C275-kutu}



\begin{quote}
Before the new one: say the Tamil for a feather, then the Tamil for a tail.
\end{quote}

\subsection*{You'll want to know: \ta{கூடு}}
\textbf{\ta{கூடு}} — \emph{kūṭu} — \textquotedblleft{}a nest\textquotedblright{}.

\begin{grammarlens}[title={What you've built}]
One more word for this chapter.
\end{grammarlens}

\subsection*{Your turn}
\begin{itemize}
\raggedright
  \item \emph{Say it:} \emph{kūṭu}
  \item \emph{Say it:} \emph{kūṭu}, once more
  \item \emph{Say it:} say \emph{kūṭu}
  \item \emph{Recall:} say the Tamil for a feather, then the Tamil for a tail, then say \emph{kūṭu} again
\end{itemize}

\begin{tcolorbox}[breakable,colback=teal!4,colframe=teal!35!black,title={Before you move on}]
What does \ta{கூடு} mean? (\textquotedblleft{}A nest\textquotedblright{}.) Say it once more.
\end{tcolorbox}
\section[\texorpdfstring{\ta{செடி} (ceṭi)}{செடி (ceṭi)}]{\ta{செடி} (ceṭi) — a plant}

\label{lesson:TA-C275-ceti}



\begin{quote}
Before the new one: say the Tamil for a tail, then the Tamil for a nest.
\end{quote}

\subsection*{You'll want to know: \ta{செடி}}
\textbf{\ta{செடி}} — \emph{ceṭi} — \textquotedblleft{}a plant\textquotedblright{}.

\begin{grammarlens}[title={What you've built}]
One more word for this chapter.
\end{grammarlens}

\subsection*{Your turn}
\begin{itemize}
\raggedright
  \item \emph{Say it:} \emph{ceṭi}
  \item \emph{Say it:} \emph{ceṭi}, once more
  \item \emph{Say it:} say \emph{ceṭi}
  \item \emph{Recall:} say the Tamil for a tail, then the Tamil for a nest, then say \emph{ceṭi} again
\end{itemize}

\begin{tcolorbox}[breakable,colback=teal!4,colframe=teal!35!black,title={Before you move on}]
What does \ta{செடி} mean? (\textquotedblleft{}A plant\textquotedblright{}.) Say it once more.
\end{tcolorbox}
\section[\texorpdfstring{\ta{கொடி} (koṭi)}{கொடி (koṭi)}]{\ta{கொடி} (koṭi) — a creeper, a vine; a flag}

\label{lesson:TA-C275-koti2}



\begin{quote}
Before the new one: say the Tamil for a nest, then the Tamil for a plant.
\end{quote}

\subsection*{You'll want to know: \ta{கொடி}}
\textbf{\ta{கொடி}} — \emph{koṭi} — \textquotedblleft{}a creeper, a vine; a flag\textquotedblright{}.

\begin{grammarlens}[title={What you've built}]
One more word for this chapter.
\end{grammarlens}

\subsection*{Your turn}
\begin{itemize}
\raggedright
  \item \emph{Say it:} \emph{koṭi}
  \item \emph{Say it:} \emph{koṭi}, once more
  \item \emph{Say it:} say \emph{koṭi}
  \item \emph{Recall:} say the Tamil for a nest, then the Tamil for a plant, then say \emph{koṭi} again
\end{itemize}

\begin{tcolorbox}[breakable,colback=teal!4,colframe=teal!35!black,title={Before you move on}]
What does \ta{கொடி} mean? (\textquotedblleft{}A creeper, a vine; a flag\textquotedblright{}.) Say it once more.
\end{tcolorbox}
